\input{../preamble.tex}
\SetDate[22/05/2026]
\setbeamertemplate{page number in head/foot}{}

\usecolortheme{crane}

\begin{document}

\section{Automatismes}

\begin{frame}

\centering \huge
Automatismes

\large
Calculatrice interdite

\end{frame}

\subsection{Signe d'image}

\framedelayed{
	\begin{center}
	\begin{tikzpicture}[>=stealth, scale=1]
	\begin{axis}[xmin = -3, xmax=3, ymin=-2.5, ymax=3.5, y=15pt, x=40pt, axis x line=middle, axis y line=middle, axis line style=->, xlabel={}, ylabel={}, xtick = {-4, -3, ..., 4}, ytick = {-4, -3, ..., 4}, grid=both, clip=false]
		
		% (f)
		\addplot[RED_E, very thick, domain =-3:3, samples=50] {6*exp(-x^2/5)-2.7}  node[right] {$\C_f$};
		
		% (g)
		\addplot[GREEN_E, very thick, domain =-3:3, samples=50] {(x+2.5)*x*(x-3)/5}  node[right] {$\C_g$};
	\end{axis}
	\end{tikzpicture}
	\end{center}
	\boxAB{
		Donner le signe de $f(-2)$.
	}{
		Donner le signe de $g(1,414)$.
	}
	Réponse attendue : strictement positif, négatif, ou nul.
}{
	\begin{center}
	\begin{tikzpicture}[>=stealth, scale=1]
	\begin{axis}[xmin = -3, xmax=3, ymin=-2.5, ymax=3.5, y=15pt, x=40pt, axis x line=middle, axis y line=middle, axis line style=->, xlabel={}, ylabel={}, xtick = {-4, -3, ..., 4}, ytick = {-4, -3, ..., 4}, grid=both, clip=false]
		
		% (f)
		\addplot[RED_E, very thick, domain =-3:3, samples=50] {6*exp(-x^2/5)-2.7}  node[right] {$\C_f$};
		
		% (g)
		\addplot[GREEN_E, very thick, domain =-3:3, samples=50] {(x+2.5)*x*(x-3)/5}  node[right] {$\C_g$};
	\end{axis}
	\end{tikzpicture}
	\end{center}
	\boxAB{
		$f(-2) = 0$.
	}{
		$g(1,414) < 0$.
	}
}

\framedelayed{
	\begin{center}
	\begin{tikzpicture}[>=stealth, scale=1]
	\begin{axis}[xmin = -2.5, xmax=2.3, ymin=-3.8, ymax=3.5, y=15pt, x=50pt, axis x line=middle, axis y line=middle, axis line style=->, xlabel={}, ylabel={}, xtick = {-4, -3, ..., 4}, ytick = {-4, -3, ..., 4}, grid=both, clip=false]
		
		% (f)
		\addplot[no marks, RED_E, -,  very thick] expression[domain=-2.5:2.3, samples=101]{1/(x-2.7) + 1/(x+2.7) + 2*x}
		node[right]{$\mathcal{C}_f$};
		
		% (g)
		\addplot[no marks, GREEN_E, -, very thick] expression[domain=-2.5:2.3, samples=101]{(x-1)*(x-2)*(x+2)/2}
		node[right]{$\mathcal{C}_g$};
	\end{axis}
	\end{tikzpicture}
	\end{center}
	\boxAB{
		Donner le signe de $f(1,618)$.
	}{
		Donner le signe de $g(1)$.
	}
}{
	\begin{center}
	\begin{tikzpicture}[>=stealth, scale=1]
	\begin{axis}[xmin = -2.5, xmax=2.3, ymin=-3.8, ymax=3.5, y=15pt, x=50pt, axis x line=middle, axis y line=middle, axis line style=->, xlabel={}, ylabel={}, xtick = {-4, -3, ..., 4}, ytick = {-4, -3, ..., 4}, grid=both, clip=false]
		
		% (f)
		\addplot[no marks, RED_E, -,  very thick] expression[domain=-2.5:2.3, samples=101]{1/(x-2.7) + 1/(x+2.7) + 2*x}
		node[right]{$\mathcal{C}_f$};
		
		% (g)
		\addplot[no marks, GREEN_E, -, very thick] expression[domain=-2.5:2.3, samples=101]{(x-1)*(x-2)*(x+2)/2}
		node[right]{$\mathcal{C}_g$};
	\end{axis}
	\end{tikzpicture}
	\end{center}
	\boxAB{
		$f(1,618) > 0$.
	}{
		$g(1) = 0$.
	}
}


\subsection{Signe de produit}

\framedelayed{
	Le nombre $x$ est strictement positif et le nombre $y$ est strictement négatif.
	\boxAB{
		Donner le signe de $x y^2$.
	}{
		Donner le signe de $\dfrac{x^2}{y}$.
	}
}{
	Le nombre $x$ est strictement positif et le nombre $y$ est strictement négatif.
	\boxAB{
		$x y^2 > 0$ (positif par positif).
	}{
		$\dfrac{x^2}{y} < 0$ (positif par négatif).
	}
}


\subsection{Inéquation}

\framedelayed{
	Quels nombres réels $x$ vérifient l'inégalité suivante ? 
	Un intervalle est attendu.
	\boxAB{
		$2x-6 \geq 0$ 
	}{
		$3x - 3 \geq 0$ 
	}
}{
	\boxAB{
		$2x-6 \geq 0 \stackrel{+6}{\iff} 2x \geq 6 \stackrel{\times\frac12}{\iff} x \geq 3$ 
		\\
		$x \in [3 ; \pinfty[$
	}{
		$3x - 3 \geq 0 \stackrel{+3}{\iff} 3x \geq 3 \stackrel{\times\frac13}{\iff} x \geq 1$
		\\
		$x \in [1 ; \pinfty[$.
	}
}



\section{Solutions}

\shipoutAnswer

\end{document}